\documentclass[10pt,a4paper]{amsart}
\usepackage[margin=19mm]{geometry}
\usepackage{amsmath,amssymb,mathtools}
\usepackage[scr=rsfs,cal=euler]{mathalfa}
\usepackage{xfrac}
\usepackage[unicode,hidelinks,bookmarks=false]{hyperref}
\newcommand{\ie}{i.e.\ }
\newcommand{\cf}{cf.\ }
\newcommand{\resp}{respectively\ }
\newcommand{\Subsection}{\subsection}
\providecommand{\nobreakdash}{\nobreak\hbox{-}}
\setlength{\emergencystretch}{2em}
\multlinegap0pt
\usepackage{amssymb}
\usepackage{float}
\usepackage{xcolor}
\newtheorem{theorem}{Theorem}
\newtheorem{lemma}{Lemma}
\newtheorem{proposition}{Proposition}
\newcommand{\sym}{\operatorname{sym}}
\title{Universal murmuration and Hecke augmentation}
\author{Shenghao Hua, Chung Pang Mok}
\date{Source: arXiv:2609.00777v1 (CC BY 4.0)}
\begin{document}
\maketitle

\noindent\emph{Source and licence.} Universal murmuration and Hecke augmentation. Version arXiv:2609.00777v1 (CC BY 4.0), distributed by arXiv under the Creative Commons Attribution 4.0 International licence, http://creativecommons.org/licenses/by/4.0/, as recorded in the arXiv licence metadata for this exact version and shown on the abstract page https://arxiv.org/abs/2609.00777v1. Full licence notice: \texttt{licenses/arxiv-2609.00777-CC-BY-4.0.txt}.\par\smallskip

\noindent\textit{Editorial scope.} Counted unit: the article's proposition prop:hyperbolic-one-r (source0.tex lines 1779--1806). The article's complete original proof of it is delivered with it (source0.tex lines 1808--1968, one \texttt{proof} environment of 161 lines). No mathematics is written, repaired or re-derived: the statement and the proof are the article's own lines, in the article's own notation, and no retained line is substituted or re-flowed.  Retained uncounted as the source-local context the proof needs: lines 835--897; lines 1165--1181; lines 1220--1259; lines 1563--1584; lines 1686--1697.  Nothing was omitted from any retained span, so the package carries no omission record.  No retained line contains a citation, so the wrapper carries no bibliography and no bibliography entry.  No retained line contains a figure environment, an image-inclusion command or drawing code, so no image is delivered and the package declares no asset dependency.  Editorial formatting only, in addition: an AMS \texttt{amsart} preamble that restates the article's own declarations and macros used by the retained lines, namely the environments \texttt{theorem}, \texttt{lemma}, \texttt{proposition} on separate counters, together with the standard packages those lines need.  The retained environments print in this document as Theorem 1, Lemma 1, Proposition 1 in order of appearance, whereas the article prints the same items with the numbering of its own full document.  The complete native source of the article as distributed in the arXiv e-print for this exact version is bundled unchanged as \texttt{sources/209061/source0.tex}.
\begin{theorem}\label{thm:scaled-murmuration}
Let $k\geq 2$ be fixed.  Let $P\neq 2$ be a prime, and suppose that $P$, $X$, and $Y$ tend to infinity.
Assume that
\[
Y=(1+o(1))X^{1-\delta_1},
\qquad
P^k\ll X^{1+\delta_2},
\]
where
\[
0<2\delta_2<\delta_1<\frac18,
\qquad
(k-1)\delta_2<1.
\]
Put
\[
\eta=\frac{\delta_1}{2}-\delta_2,
\qquad
y=\frac{P^k}{X}.
\]
Let
\[
A=\prod_p\left(1+\frac{p}{(p+1)^2(p-1)}\right),
\]
\[
B=\prod_p\left(1-\frac{p}{(p^2-1)^2}\right),
\]
and
\[
C(r)=\prod_{p\mid r}
\left(1+\frac{p^2}{p^4-2p^2-p+1}\right).
\]
Then, for every $\varepsilon>0$,
\begin{multline*}
\frac{
\displaystyle
\sideset{}{'}\sum_{N\in[X,X+Y]}
\ \sum_{f\in H_{\mathrm{new}}(N)}
P^{k/2}\lambda_{\sym^k f}(P)\,\varepsilon(f)
}{
\displaystyle
\sideset{}{'}\sum_{N\in[X,X+Y]}
\ \sum_{f\in H_{\mathrm{new}}(N)}1
}
\\
=
\frac{12}{\pi D_0}
\left(
A\sqrt y
+
B\sum_{1\leq r\leq 2\sqrt y}
C(r)\sqrt{4y-r^2}
-\pi y
\right)
+
O_{k,\varepsilon}\left(
X^{-\eta+\varepsilon}
+X^\varepsilon\frac{\sqrt y}{P}
+\frac{1+y}{P}
\right).
\end{multline*}
Here the prime indicates that the sum is restricted to squarefree integers $N$.
\end{theorem}

\begin{lemma}\label{lem:p-coprime}
Assume the hypotheses of Theorem~\ref{thm:scaled-murmuration}.
If
\[
0<r\leq\frac{2P^{k/2}}{\sqrt X}
\]
and
\[
d^2\mid r^2N-4P^k,
\qquad
P\nmid N,
\]
then, for sufficiently large $X$,
\[
P\nmid rd.
\]
\end{lemma}

\begin{lemma}\label{lem:admissible-classes}
Suppose that $P\nmid rd$.  There is a set
$\mathcal R_{r,d}(Q)\subset\mathbb Z/d^2\mathbb Z$ such that
\[
\mathcal A_{r,d}(Q)
=
\left\{
N\in\mathbb Z_{>0}:
\mu^2(N)=1,\ P\nmid N,\
N\bmod d^2\in\mathcal R_{r,d}(Q)
\right\}.
\]
Every class in $\mathcal R_{r,d}(Q)$ is coprime to $d$, and
\[
\#\mathcal R_{r,d}(Q)
=
\begin{cases}
1,&(r,d)=1,\quad 2\nmid rd,\\
1,&(r,d)=1,\quad 2\mid r,\\
1,&(r,d)=2,\quad 2\Vert d,\\
2,&(r,d)=2,\quad 4\mid d,\\
0,&\text{otherwise}.
\end{cases}
\]
If $d$ is even, every integer represented by an admissible class is odd.
Writing
\[
s=\frac{N(r^2N-4Q)}{d^2},
\]
the parity information is
\[
\begin{array}{c|c}
\text{case}&\text{parity of }s\\ \hline
2\Vert d,\ 2\Vert r&\text{even}\\
2\Vert d,\ 4\mid r&\text{odd}\\
4\mid d&\text{one even class and one odd class}.
\end{array}
\]
In particular, the number and parity type of these classes are independent of $k$.
\end{lemma}

\begin{lemma}\label{lem:small-n}
For $0\leq\sigma,\tau\leq1$,
\begin{equation*}
\begin{aligned}
\sum_{\substack{n\leq Y^\sigma\\d\leq Y^\tau}}
\frac{\mathcal S_{d,n,r}}{nd}
&=
\frac{Y}{\zeta(2)}BC(r)
+O\left(\frac{Y}{P^2}\right)\\
&\quad+
O_\varepsilon\left(
(QXY)^\varepsilon
\left(
\sqrt X\,Y^{\sigma/2}
+Y^{\tau+3\sigma/2}
+Y^{1-2\tau}
+Y^{1-\sigma/5}
\right)
\right).
\end{aligned}
\end{equation*}
\end{lemma}

\begin{lemma}\label{lem:large-n}
Let $T\geq Y^\sigma$.  For $0\leq\sigma,\tau\leq1$,
\[
\sum_{\substack{Y^\sigma<n\leq T\\d\leq Y^\tau}}
\frac{\mathcal S_{d,n,r}}{nd}
\ll_\varepsilon
(QXT)^\varepsilon
\left(
\sqrt X+Y^{1-\sigma/2}+\sqrt Y\,T^{1/4}
\right).
\]
\end{lemma}

\begin{proposition}\label{prop:hyperbolic-one-r}
Uniformly in this range of $r$, one has
\begin{multline*}
\mathcal H_k(r)
=
\frac{Y\sqrt{4P^kX-r^2X^2}}{\zeta(2)\pi}BC(r)
\\
+
O_{k,\varepsilon}\left(
(P^kXY)^\varepsilon
\left(
(YP^kX)^{3/5}
+
\frac{Y^2P^{k/2}}{\sqrt X}
+
r\sqrt X\,Y^{3/2}
+
P^{k/2}XY^{5/18}
\right.\right.
\\ \left.\left.
+
P^{k/2}\sqrt X\,Y^{8/9}
\right)
\right)
+
O_k\left(P^{k/2-2}\sqrt X\,Y\right).
\end{multline*}
\end{proposition}

\begin{proof}
Put
\[
\Delta_r=4QX-r^2X^2.
\]
For $N$ in the interval, the class number formula and
Lemma~\ref{lem:admissible-classes} give
\begin{equation}\label{eq:class-number-double-sum}
H_1(r^2N^2-4QN)
=
\frac1\pi
\sum_{\substack{d^2\mid r^2N-4Q\\
N\bmod d^2\in\mathcal R_{r,d}(Q)}}
\frac{\sqrt{4QN-r^2N^2}}{d}
L\left(1,\chi_{\frac{N(r^2N-4Q)}{d^2}}\right)
+
O(1).
\end{equation}
To justify that this accounts for every square divisor in the Hurwitz class number, note that $N$ is squarefree and $P\nmid N$.
A square factor which uses a prime from both $N$ and $r^2N-4Q$ can therefore occur only at $2$.
If $N=2N_0$, with $N_0$ odd, removal of this extra factor $4$ would leave
\[
\frac{r^2N_0^2-2QN_0}{d^2}.
\]
Here $d$ is odd: the odd parts of $N_0$ and
$r^2N_0-2Q$ are coprime, and the latter has $2$-adic valuation at most $1$.
This is $2\pmod4$ when $r$ is even and $3\pmod4$ when $r$ is odd, so its Gauss class number is zero.
Thus only the divisors displayed in~\eqref{eq:class-number-double-sum} contribute.
The discriminant condition is encoded by Lemma~\ref{lem:admissible-classes}.
Moreover $d<2\sqrt Q$, and Lemma~\ref{lem:p-coprime} gives $P\nmid rd$.
The $O(1)$-term in~\eqref{eq:class-number-double-sum}, summed over $N$, is $O(Y)$.
Since the present range is nonempty only when $Q>(X+Y)/4$, this is absorbed by $(QXY)^{3/5}$.

Let $0<\tau<1/2$.
The bound
\[
L(1,\chi)\ll_\varepsilon(QX)^\varepsilon
\]
and the fact that each $d$ gives at most two residue classes show that the part of~\eqref{eq:class-number-double-sum} with $d>Y^\tau$, summed over $N$, is
\begin{align}
\ll_\varepsilon
(QX)^{1/2+\varepsilon}
\sum_{Y^\tau<d<2\sqrt Q}
\left(\frac{Y}{d^2}+1\right)\frac1d
\ll_\varepsilon
(QX)^{1/2+\varepsilon}
\left(Y^{1-2\tau}+\log(2Q)\right).
\label{eq:large-d}
\end{align}

For $d\leq Y^\tau$, we freeze the square-root factor at $N=X$.
Uniformly in the interior range,
\[
\left|
\sqrt{4QN-r^2N^2}-\sqrt{\Delta_r}
\right|
\ll
\frac{\sqrt Q\,Y}{\sqrt X}
+r\sqrt{XY}.
\]
After summing over $N$ and $d$, the resulting error is
\begin{equation}\label{eq:square-root-variation}
O_\varepsilon\left(
(QXY)^\varepsilon
\left(
\frac{\sqrt Q\,Y^2}{\sqrt X}
+r\sqrt X\,Y^{3/2}
\right)
\right).
\end{equation}

For a quadratic character of discriminant $D$, partial summation together with P\'olya--Vinogradov gives
\[
L(1,\chi_D)
=
\sum_{n\leq T}\frac{\chi_D(n)}n
+
O\left(\frac{\sqrt{|D|}\log(2|D|)}T\right).
\]
Since $|D|\ll QX/d^2$, inserting this in the remaining part of~\eqref{eq:class-number-double-sum} produces the truncation error
\begin{equation}\label{eq:L-truncation-error}
O_\varepsilon\left(
(QXY)^\varepsilon\frac{QXY}{T}
\right).
\end{equation}
The identity
\[
\left(\frac{N(r^2N-4Q)/d^2}{n}\right)
=
\left(\frac Nn\right)
\left(\frac{(r^2N-4Q)/d^2}{n}\right)
\]
now yields
\begin{multline}\label{eq:hyperbolic-master}
\mathcal H_k(r)
=
\frac{\sqrt{\Delta_r}}{\pi}
\sum_{\substack{d\leq Y^\tau\\n\leq T}}
\frac{\mathcal S_{d,n,r}}{nd}
\\
+
O_\varepsilon\left(
(QXY)^\varepsilon
\left(
\frac{QXY}{T}
+\frac{\sqrt Q\,Y^2}{\sqrt X}
+r\sqrt X\,Y^{3/2}
+\sqrt{QX}\,Y^{1-2\tau}
\right)
\right),
\end{multline}
where the logarithm in~\eqref{eq:large-d} has been absorbed into the last term.
Indeed, $\log(2Q)\ll\log X=o(Y^{1-2\tau})$ because $\tau<1/2$ and the short-window hypotheses make $Y$ a fixed positive power of $X$.

Choose
\[
\tau=\frac1{18},
\qquad
\sigma=\frac59.
\]
Lemmas~\ref{lem:small-n} and~\ref{lem:large-n} give
\begin{equation}\label{eq:arithmetic-double-sum}
\sum_{\substack{d\leq Y^\tau\\n\leq T}}
\frac{\mathcal S_{d,n,r}}{nd}
=
\frac{Y}{\zeta(2)}BC(r)
+
O\left(\frac{Y}{P^2}\right)
+
O_\varepsilon\left(
(QXT)^\varepsilon
\left(
\sqrt X\,Y^{5/18}
+Y^{8/9}
+\sqrt Y\,T^{1/4}
\right)
\right).
\end{equation}
The range of $r$ can be nonempty only if $4Q>X+Y$.
Hence, with
\[
T=(YQX)^{2/5},
\]
\[
\frac TY
=
\frac{(YQX)^{2/5}}Y
\gg X^{(1+3\delta_1)/5+o(1)}.
\]
Thus, in particular, $T>Y$, and Lemma~\ref{lem:large-n} applies.
Finally,
\[
\frac{QXY}{T}
\asymp
\sqrt{QXY}\,T^{1/4}
\asymp
(QXY)^{3/5}.
\]
Substituting~\eqref{eq:arithmetic-double-sum} into \eqref{eq:hyperbolic-master}, using $\sqrt{\Delta_r}\leq\sqrt{QX}$, and recalling $Q=P^k$, gives exactly the error stated in the
proposition.
\end{proof}

\end{document}
